\documentclass[10pt,journal]{IEEEtran}
\usepackage{microtype}
\usepackage{amsmath,amssymb,mathtools}
\usepackage{booktabs,multirow,array}
\usepackage{graphicx}
\usepackage{cite}
\usepackage{xcolor}
\usepackage{url}
\emergencystretch=1em
% Retain IEEEtran's journal baseline; page length is controlled by prose and
% float placement rather than compressed line spacing.
\setlength{\textfloatsep}{5pt plus 1pt minus 1pt}
\setlength{\floatsep}{4pt plus 1pt minus 1pt}
\setlength{\intextsep}{5pt plus 1pt minus 1pt}
\setlength{\abovedisplayskip}{3pt plus 1pt minus 1pt}
\setlength{\belowdisplayskip}{3pt plus 1pt minus 1pt}
\renewcommand{\arraystretch}{0.98}
\newcommand{\RiskSafe}{risk-constrained verifier}
\newcommand{\TailRisk}{tail-risk counterfactual}
\newcommand{\PaymentSafe}{payment-certified verifier}
\newcommand{\dt}{\Delta t}

\title{Security-Aware Workload-State Counterfactual Verification and\\
Nodal Net-Value Settlement for Spatially Coupled AI Data Centers}
\author{Anonymous Authors}

\begin{document}
\maketitle

\begin{abstract}
Artificial-intelligence (AI) data centers are flexible grid resources, but a
meter-only demand-response baseline cannot show that credited work was
committed before the event, nor can it distinguish a local reduction from
spatial rebound. This paper develops an auditable submit--execute--settle
verifier that carries one declared workload state from counterfactual
scheduling to signed nodal settlement. A sparse cumulative linear program and
an exact job-indexed witness enforce release, deadline, energy, graphics-
processing-unit (GPU) nameplate, and regional-capacity constraints. The
reported risk profile is a validation-fitted simplex under total and daily
conditional-value-at-risk budgets; a separate submission-masked gate policy
and a fixed/flexible power-conversion certificate keep execution telemetry out
of contract formation. On the 121-day trace inventory, the locked
mechanism-isolation panel gives normalized root-mean-square error (nRMSE)
0.314, false credit 1.809 MWh/day, and $F_1$ score 0.470, compared with
0.339, 1.823 MWh/day, and 0.418 for the selected single projection. The ledger contains 71,144 submissions and 71,128 matched
executions; its indexed witness has 12,296,675 job--slot variables and a
declared-energy residual below $10^{-15}$ MWh. A 54-day gate replay uses only
pre-event declarations, while an independent network replay evaluates the
same witness under all 37 finite RTS-24 contingencies and two homogeneous
scale transforms. Fixed facility demand is carried separately, so the raw
$q_{99}$ endpoint remains an auditable activation scenario rather than a
silently clipped stress case. The certificate reports workload feasibility,
information provenance, and signed N--1 value without treating public traces
as co-located utility measurements.
\end{abstract}

\noindent\textbf{Index Terms---}AI data centers, counterfactual verification,
demand response, security-constrained dispatch, nodal settlement.

\section{Introduction}

The growth of large-model training and inference is converting data centers
from nearly static industrial loads into large, software-defined grid
resources. Interactive inference is latency constrained, elastic inference can
be delayed briefly, and batch graphics-processing-unit (GPU) work can be deferred for hours or migrated
between regions. These capabilities enable demand response (DR), but they also
invalidate the usual interpretation of a meter decrease. A reduction at one
facility can be followed by delayed consumption or simultaneous rebound at
another node. Paying the positive local decrease can therefore reward a
redistribution that creates little or negative network value.
Public traces already expose multi-week bursty workload arrivals and tens of
thousands of scheduler jobs, while facility telemetry reports GPU energy at a
different measurement boundary \cite{wang2024burstgpt,samsi2021supercloud}.

Conventional DR measurement and verification estimates the unobserved
no-event demand from historical meter data. Arithmetic historical baselines
remain common in market practice under published market guidance
\cite{caiso2017baseline}, despite known bias and manipulation incentives
\cite{wang2022baseline}. Regression-based facility models quantify
baseline uncertainty but do not encode a computing workload ledger
\cite{wang2022baseline}. More broadly, DR
research establishes the economic role of responsive load, yet a meter-only
estimator cannot determine whether queued jobs could actually have produced
its predicted trajectory.

Data-center scheduling studies provide the missing physical state. Temporal
deferral, dynamic capacity, and geographical load balancing can reduce energy
cost and peak demand. Recent work assesses data-driven flexibility from
production traces and system-level
security-constrained coordination, while balancing-market and clean-computing
studies optimize spatial--temporal data-center dispatch
\cite{cao2022flexibility,cao2024nonwire}. A recent power-system review also
identifies quality-of-service (QoS), service-level-agreement (SLA), and
market-participation constraints that are absent from meter-only accounting;
clean-computing coordination provides a complementary dispatch context
\cite{takci2025flexibility,riepin2025clean}.
Hierarchical optimal-consumption models further connect facility efficiency
with response quality.
Virtual-link market clearing and space--time remuneration establish how
geographical and temporal transfers can enter electricity markets
\cite{zhang2020virtuallinks,zhang2022remunerating}. Receding-horizon operation
and uncertainty-aware frequency-regulation coordination further establish the
need to retain future arrivals and multi-site state
\cite{zhang2023receding}. These methods optimize
operator-controlled schedules; verification of an independently submitted
DR counterfactual against an auditable workload ledger remains unresolved. This separation between
computing feasibility and market verification motivates the auditable ledger
used here. Recent production-trace studies quantify data-center flexibility and
load-aggregator coordination, but leave this ledger-to-settlement coupling open
\cite{dcaopt2024,caprara2026}. In particular, cross-regional dispatchable-capacity
evaluation and coupled-regulation modeling quantify operational flexibility but
do not provide a provenance-bound settlement certificate
\cite{han2026dispatchable,chen2021idccoupling}. Recent whole-facility measurements likewise
require a separate bottom-up conversion from GPU power to facility demand;
the present study therefore keeps raw energy calibration, benchmark scaling,
and network placement as separate evidence layers.

The second gap is electrical. Geographic workload control is commonly valued
with energy prices or gross megawatts. Such accounting ignores the fact that
nodal marginal value changes when congestion or a contingency constraint
binds. Direct-current (DC) power flow and linear distribution factors provide transparent
network sensitivities \cite{stott2009dc,zimmerman2011matpower}; N--1 dispatch
can impose every credible post-contingency line limit through line-outage
distribution factors (LODFs) \cite{tejada2018lodf}. No existing data-center DR
verification framework jointly certifies workload feasibility, preserves
signed spatial response, and settles the resulting security-aware network
value. Alternating-current (AC) feasibility is a distinct nonlinear check,
not a consequence of the DC certificate. Public workload and facility telemetry can support an observational
replay of that ledger, but it does not constitute a causal utility event; a
rigorous study must therefore separate measured-execution replay, counterfactual
construction, and network-value verification to avoid conflating them with one
field intervention.

This paper makes three power-system contributions. 1) It defines a typed
submit--execute--settle evidence chain: an exact job-indexed counterfactual is
formed only from scheduler declarations (release, runtime, requested GPUs,
and declared energy entitlement), whereas the execution ledger is retained
for post-event scoring and the settlement ledger is generated only after the
coupling certificate passes. The indexed witness is independently rechecked
before its regional profile is valued in the N--1 model. 2) It formulates a
workload-state verifier with an explicit information boundary and a
validation-fitted total-plus-CVaR risk fit, while keeping the pointwise
false-credit cap in a separately auditable payment profile. A separately
parameterized event policy and a binding-capacity stress cohort test whether
the declaration-based feasible set remains nonempty when regional limits,
service floors, and declared windows are jointly active. 3) It develops signed
space--time nodal remuneration and a payment-activation rule over held-out
power-conversion factors: the raw $q_{99}$ endpoint is retained, but a payment
can activate only when the precommitted fixed/flexible capacity certificate
covers it. The central result is a
dimension-preserving coupling invariant that carries one declared workload
state into physical verification, network valuation, and settlement; the
individual linear-program (LP), dispatch, and allocation primitives are standard and are not
claimed as new in isolation. Because public traces do not identify utility
geography, spatial conclusions are conditional on the declared feature-
stratified trace scenario and the complete permutation/penetration panel.
The certificate assumes trusted scheduler and telemetry provenance: a digest
detects post-commitment record changes but cannot authenticate semantically
false fields.

The remainder of the paper is organized as follows. Section~\ref{sec:model}
defines the spatial workload and network settlement problem. Section~\ref{sec:method}
presents the verifier, security-aware value rule, and allocation mechanism.
Section~\ref{sec:design} specifies the data, baselines, and locked protocol.
Section~\ref{sec:results} reports the core counterfactual, execution-replay,
job-level, security, payment, continuous-horizon, and provenance experiments.
Section~\ref{sec:conclusion} concludes with the deployment boundary and future
validation requirements.

\section{System Model and Problem Formulation}\label{sec:model}

\subsection{Spatial workload state}

Let $\mathcal S$, $\mathcal K$, $\mathcal D$, and $\mathcal T$ denote source
regions, workload classes, destination data centers, and 15-minute intervals;
the symbol $h$ is reserved for a calendar day in statistical panels.
The measured ledger supplies arrival energy $a_{skt}$ for source $s$, class
$k$, and release interval $t$. The decision $x_{skd\tau}\geq0$ is the energy
served at destination $d$ in interval $\tau$. Its cumulative state is
\begin{equation}
 y_{skt}=\sum_{\tau=0}^{t}\sum_{d\in\mathcal D}x_{skd\tau}.
 \label{eq:cumulative}
\end{equation}
With cumulative arrivals $A_{skt}=\sum_{\tau=0}^{t}a_{sk\tau}$ and deadline
$D_k$, release and deadline feasibility are
\begin{align}
 y_{skt}&\leq A_{skt}, \label{eq:release}\\
 y_{skt}&\geq A_{sk,t-D_k},\quad t\geq D_k. \label{eq:deadline}
\end{align}
The state recursion equals service, terminal service equals total arrivals, and
$\sum_{s,k}x_{skdt}\leq\bar P_d\dt$. Facility demand is
\begin{equation}
 p_{dt}=P_d^{\rm fix}+\dt^{-1}\sum_{s,k}x_{skdt}.
 \label{eq:power}
\end{equation}
For rolling window end $T^+$ and final event-day release $T^0$, the terminal
state is not forced to consume unexpired future work. Instead,
\begin{equation}
y_{skT^+}=A_{sk,\max\{T^0,T^+-D_k\}},
\label{eq:rollingterminal}
\end{equation}
which completes every job due inside the window and every event-day arrival,
while retaining later real arrivals in the release constraints. The window
starts $D_{\max}+96$ intervals before the event day and ends $D_{\max}$
intervals after it; omitted earlier work must therefore expire before the
evaluated day.
Equations (\ref{eq:release}) and (\ref{eq:deadline}) prevent early execution
and late completion, while (\ref{eq:power}) maps the feasible service state to
the facility meter.
The aggregate variable \(x\) is a preemptive envelope for checkpointable batch
work; contiguous job intervals are retained separately for the MIT replay and
checked in Experiment~14. This two-level construction keeps the counterfactual
optimization sparse while making the task-level scope explicit.
For the task-level check, \(j\in\mathcal J_{\rm sub}\) has submitted release
\(r_j\), entitlement \(\widehat E_j\), scenario region \(s_j\), and requested GPU count
\(g_j^{\rm req}\). Its endpoint is \(d_j=r_j+R_j+B\), where
\(R_j=\lceil\mathrm{timelimit}_j/\dt\rceil\), \(B=96\) precommitted queue slots,
and the unlimited sentinel is mapped to \(H_{\infty}=128\) slots. Finite
declarations are retained exactly; the $0.001$-MW-per-GPU nameplate is an
audit-only precheck, and observed completion never enlarges \(d_j\). The indexed
service variable \(u_{j\tau}\) is constrained by
\begin{equation}
\begin{aligned}
 &u_{j\tau}\geq0,\qquad
   \sum_{\tau=r_j}^{d_j-1}u_{j\tau}=\widehat E_j,\\
 &u_{j\tau}\leq g_j^{\rm req}\bar p_{\rm GPU}\dt,\\
 &\sum_{j:s_j=d,\ r_j\leq\tau<d_j}u_{j\tau}
   \leq \bar P_d\dt,\\[-2pt]
 &\min_u\ \sum_{j,\tau}
   \left[w_{\rm batch}(\tau-r_j)+r^{\rm DR}\mathbf1\{\tau\in\mathcal E\}\right]u_{j\tau}.
\end{aligned}
\label{eq:jobindexed}
\end{equation}
Here \(\bar p_{\rm GPU}=0.001\) MW/GPU. Equation~(\ref{eq:jobindexed}) is a
globally solved, declaration-only schedule rather than a replay; its event
tariff is fixed before scoring. Experiment~19 covers all 71,144 submissions,
of which 71,128 have a positive-energy match used only for independent scoring.
The job-to-network certificate maps this witness without a second optimization.
With native assignment \(d=s_j\), let $\mathcal J_{skd}$ be the corresponding
ledger partition:
\begin{align}
 x^{\rm job}_{skd\tau}&=\sum_{j\in\mathcal J_{skd}:\,r_j\leq\tau<d_j}
 u_{j\tau},\\
 p^{\rm job}_{d\tau}&=P_d^{\rm fix}+\dt^{-1}X^{\rm job}_{d\tau},\qquad
 X^{\rm job}_{d\tau}=\sum_{s,k}x^{\rm job}_{skd\tau}.
 \label{eq:jobnetworkcoupling}
\end{align}
Experiment~22 reconstructs $X^{\rm job}$ from every stored $u_{j\tau}$,
rechecks the indexed primal, and verifies equality before network valuation;
the aggregate migration profile is only a matched-workload control. Workload
deferral and geographic assignment follow \cite{cao2022flexibility,cao2024nonwire};
rolling-state control follows \cite{zhang2023receding}; waiting and migration
costs are $w_k(\tau-t)x_{skd\tau}$ and $m_{sd}x_{skd\tau}$.

For provenance, the submit-time ledger and post-event execution ledger have
separate canonical digests. Each submitted record is represented by the
scheduler-only tuple and its digest in
\begin{equation}
\begin{aligned}
 \ell_j^{\rm sub}={}&(\mathrm{id}_j,t_j^{\rm submit},\mathrm{timelimit}_j,
 q_j,\kappa_j,\sigma_j),\\
 \mathcal L_{\rm sub}={}&\operatorname{sort}\{\ell_j^{\rm sub}:j\in\mathcal J_{\rm sub}\},\\
 h_{\rm sub}={}&\operatorname{SHA256}\!\left(\operatorname{serialize}(\mathcal L_{\rm sub})\right).
\end{aligned}
\label{eq:ledgerdigest}
\end{equation}
where $q_j$, $\kappa_j$, and $\sigma_j$ are requested-resource, class, and
scheduler-state fields. The execution ledger stores start/end times, DCGM
energy, and counts; its join and energy equality are checked only post-event.
Canonical ordering and the Secure Hash Standard \cite{nist2015fips1804} keep
$h_{\rm sub}$ in (\ref{eq:ledgerdigest}) independent of execution telemetry.
Experiment~16 reconciles the measured envelope with the pre-split 118-MW
nameplate. Scheduler/DCGM records are authoritative; the digest detects
changes but cannot attest semantic truth of submitted fields.

\subsection{Counterfactual, response, and false credit}

For event set $\mathcal E\subset\mathcal T$, $p^0$ is the no-event
counterfactual, $p^{1,\mathrm{sim}}$ the workload event trajectory, and
$p^{\rm obs}$ the independent MIT/DCGM replay. Because no utility event label is
available, the replay is observational only and never a causal outcome or
pre-decision input. In mechanism-isolation rows, signed rebound is retained as
\(r=p^0-p^{1,\mathrm{sim}}\).
For source $z$ in the observed or simulated replay, define
$\widehat r^{(z)}=[\widehat p^0-p^{(z)}]_+$ and
$r^{(z)}=[p^0-p^{(z)}]_+$; source-specific false credit is
\begin{equation}
 F^{(z)}=\dt\sum_{d,t\in\mathcal E}
 [\widehat r^{(z)}_{dt}-r^{(z)}_{dt}]_+,
 \qquad z\in\{\mathrm{obs},\mathrm{sim}\}.
 \label{eq:falsecredit}
\end{equation}
The simulated source defines the validation epigraph; the observed source is
recomputed only in locked replay. The identity
\(p^{(\mathrm{sim})}+r^{(\mathrm{sim})}=\max\{p^0,p^{(\mathrm{sim})}\}\) is audited
but is not a new meter trajectory. Before the event the operator freezes
\(p^{\rm con}\); after closure, payable response is
\begin{equation}
Q^{\rm pay}=\dt\sum_{d,t\in\mathcal E}
\min\!\left\{[\widehat p^0_{dt}-p^{\rm obs}_{dt}]_+,
 [p^{\rm con}_{dt}-p^{\rm obs}_{dt}]_+\right\}.
\label{eq:meter_cap}
\end{equation}
Equation (\ref{eq:meter_cap}) uses only the submitted baseline, frozen contract,
and closed meter. The trace-anchored $p^0$ is an offline diagnostic and never
enters fitting, gate decisions, workload optimization, or transfer formation.
Reported metrics include gross/payable response, oracle over/underpayment,
nRMSE, bias, precision, recall, and $F_1$.

\subsection{Strategic reference scheduling}

Consider an arithmetic mean of $H=10$ reference days, event probability $q$,
and DR price $\pi^{\rm DR}$. If reference day $j$ adds event-window service
$\delta R_j$, the expected payment change is
\begin{equation}
 \Delta\Pi=\frac{q\pi^{\rm DR}}{H}\sum_{j=1}^{H}\delta R_j
 -\sum_{j=1}^{H}\Delta C_j .
 \label{eq:profit}
\end{equation}
Let $c'_j$ be the right derivative of the minimum operating cost with respect
to an event-window minimum-service constraint on day $j$. Linear-program
sensitivity \cite{boyd2004convex} gives the following exact local condition:
\begin{equation}
\begin{aligned}
 \frac{q\pi^{\rm DR}}{H}>\min_j c'_j
 \quad\Longleftrightarrow\quad {}&
 \exists\ \text{a strictly profitable}\\[-2pt]
 &\text{local deviation}.
\end{aligned}
 \label{eq:threshold}
\end{equation}
At a degenerate optimum the dual at the current service level can be zero
even when the right-direction marginal cost is positive. Experiment~1
therefore certifies $c'_j$ with two independent right finite differences of
$5\times10^{-4}$ and $10^{-3}$ MWh and a $10^{-6}$ USD/MWh agreement
tolerance.
The historical-baseline premise is documented in
\cite{wang2022baseline}; (\ref{eq:profit})--(\ref{eq:threshold})
are this paper's workload-specific specialization.

\subsection{Network value}

Let $M$ be the data-center-to-bus incidence matrix and
$p_t=p_t^{\rm native}+Mp_t^{\rm dc}$ the nodal demand vector. For
generator-to-bus map $C_g$, generator bounds, and power-transfer distribution
factor (PTDF) $H$ with line ratings $\bar f$, the
piecewise-linear security-constrained economic dispatch (SCED) value is
\begin{align}
 V(p)=\min_g\;&\sum_i C_i(g_i) \label{eq:sced}\\
 {\rm s.t.}\;&\mathbf1^\top g=\mathbf1^\top p,\quad
 \underline g\leq g\leq\bar g,\nonumber\\
 &-\bar f\leq H(C_gg-p)\leq\bar f.\nonumber
\end{align}
The incidence mapping aggregates device demand at its connection bus, and
(\ref{eq:sced}) is the standard lossless DC representation
\cite{stott2009dc,zimmerman2011matpower}. The realized signed grid value is
\begin{equation}
 \Delta V=V(p^0)-V(p^1). \label{eq:value}
\end{equation}
Positive $\Delta V$ is a credit and negative $\Delta V$ a debit; no
positive-only truncation is applied.

\section{Proposed Methodology}\label{sec:method}

\subsection{Model-based framework and matched-information pre-estimation}

The state, power, and network equations in Section~\ref{sec:model} are the
first stage of the methodology: they define the feasible workload state before
any statistical target or payment is formed. Fig.~\ref{fig:framework} shows the
complete chain. Meter history and the
submitted-job ledger enter a statistical pre-estimator and the workload
feasible set in parallel. The pre-estimator never receives execution energy
used as scoring truth. For each event day, it uses only the information
available under one of two declared protocols: lagged meter/calendar features
for online baselines, or the complete submitted ledger for equally informed
post-event comparisons. Ridge regression and gradient boosting are used as
standard matched-information comparators \cite{hoerl1970ridge,friedman2001gradient};
extremely randomized trees and synthetic control provide two additional
nonlinear and donor-pool comparators \cite{geurts2006extratrees,abadie2010synthetic}.
Two stronger comparators are also solved: median-loss gradient boosting with
the complete submitted ledger and a nonnegative synthetic control whose
weights globally minimize pre-event absolute error over 28 donor days.
To separate the effect of projection from that of pre-estimation, the
complete-ledger quantile output is also passed through the same exact
workload-feasible projection, with its penalty chosen by contiguous
validation folds.
The same scheduler declaration view is supplied to the post-event learner, the
selected single projection, and the proposed convex verifier; execution energy
and completion timestamps are withheld until the locked scoring pass.

\begin{figure*}[t]
\centering
\includegraphics[width=0.82\textwidth]{figures/fig0_framework.pdf}
\caption{Verification-to-settlement framework. Validation labels select the
counterfactual risk contract; locked execution truth is used only for scoring.
The right branch computes aggregate N--1 value and exact participant shares.}
\label{fig:framework}
\end{figure*}

\subsection{Sparse cumulative workload-state verifier}

For statistical estimate $\widetilde p$ and each predeclared penalty $\rho_\ell$,
the first-stage program solves
\begin{equation}
 \min_{x,y,z}\ \mathcal C_{\rm w}(x)+\rho_\ell\sum_{d,t}z_{dt}
 \quad {\rm s.t.}\quad
 z_{dt}\geq\pm(p_{dt}(x)-\widetilde p_{dt}),
 \label{eq:projection}
\end{equation}
subject to (\ref{eq:cumulative})--(\ref{eq:power}). The absolute-value
epigraph is standard \cite{boyd2004convex}; its use to project a meter
baseline onto the auditable workload set is the study-specific step. Six
globally solved schedules $p^{(\ell)}$ share all arrivals and constraints, so
their validation-only simplex combination
$\bar p=\sum_\ell\alpha_\ell p^{(\ell)}$ remains feasible.

The coefficients minimize validation squared error subject to total and
daily-tail mechanism-isolation false-credit budgets; conditional value at risk
(CVaR) measures the latter \cite{rockafellar2000cvar}. Let
$\ell^\star$ be the single candidate with minimum contiguous-fold nRMSE and
$B_j^{\rm cap}=F_j^{(\mathrm{sim})}(\ell^\star)$. The observed-meter source is
audited only after the locked decision. For reserve $\beta\in(0,1]$,
\begin{align}
 \sum_jF_j^{(\mathrm{sim})}(\alpha)&\leq\beta\sum_jB_j^{\rm cap},\label{eq:totalrisk}\\
 {\rm CVaR}_{0.75}\!\left(F_j^{(\mathrm{sim})}(\alpha)\right)
 &\leq\beta{\rm CVaR}_{0.75}(B_j^{\rm cap}).\label{eq:cvar}
\end{align}
Implementation linearizes the positive part with $f_n\geq0$, a CVaR threshold
$\nu\geq0$, and tail slacks $u_j\geq0$ for each event sample and validation day:
 \begin{align}
  f_n&\geq p^\alpha_n-p^{(\mathrm{sim})}_n-r^{(\mathrm{sim})}_n,\\
  F_j^{(\mathrm{sim})}&=\sum_{n\in j}f_n,\\
  F_j-\nu-u_j&\leq0,\\
  \nu+\frac{1}{0.25J}\sum_j u_j
  &\leq\beta\,\mathrm{CVaR}_{0.75}(B_j^{\rm cap}).
  \label{eq:riskepigraph}
 \end{align}
Here $r^{(\mathrm{sim})}_n=[p^0_n-p^{(\mathrm{sim})}_n]_+$; the explicit simulated
and true terms remain in the model. The epigraph is never substituted for the
frozen contract profile $p^{\rm con}$. Together with the simplex and box
constraints, (\ref{eq:riskepigraph}) is a convex quadratic program. HiGHS
supplies a feasible point, after which stationarity, epigraph, simplex, and
bound residuals are recomputed and stored. Four contiguous folds select $\beta$
by validation nRMSE and then risk ratio; held-out noninferiority is an outcome,
not a hidden constraint. A separate payment-contract projection targets $\bar p$
while imposing, for every event sample,
\begin{equation}
p^{\rm safe}_{dt}\leq p^{(\ell^\star)}_{dt}. \label{eq:envelope}
\end{equation}
The pointwise cap in (\ref{eq:envelope}) is the contractual upper band, evaluated
on every event slot before settlement and never substituted for the risk profile.
\paragraph{Two-sided credit band.}
The upper envelope is paired with a lower band tied to the independently fitted
target. Define \(p^{\rm floor}_{dt}=\min\{p^{(\ell^\star)}_{dt},\bar p_{dt}\}\);
the event-slot constraints are
\begin{equation}
 p^{\rm floor}_{dt}-\varepsilon
 \leq p^{\rm safe}_{dt}\leq p^{(\ell^\star)}_{dt},
 \qquad (d,t)\in\mathcal D\times\mathcal E,
 \label{eq:twosided}
\end{equation}
where $\varepsilon=1$ MW is fixed before the locked split and imposed with the
workload constraints, not clipped afterward. The total and CVaR constraints
control aggregate and daily-tail exposure; the upper side is the deterministic
payment cap, while the lower side is an under-credit regularizer rather than a
distribution-free unseen-trace guarantee. The selected single projection proves
nonemptiness, and locked outputs report both margins and payable response.

\paragraph{Decision-time information boundary.}
The ex-post selector may use the complete ledger, but a gate deployment removes
all post-gate arrivals before target construction. The masked target is refit
with the matched-information model; a no-tariff LP freezes the baseline and
commits 5 MWh, while a second LP applies the DR price with unexpired state.
Reserved future jobs cannot be credited. Baseline error is scored only for the
submitted contract; the tariff plan is scored against the closed meter and
settled by (\ref{eq:meter_cap}).

To separate capacity planning from payment eligibility, the gate publishes a
causal reserve envelope. For candidate quantile $\eta$ and event day $d$,
\begin{equation}
 \widehat a^{\rm res}_{sk\tau}(d;\eta)=Q_{\eta}\!\left(\left\{a_{sk\tau}(h):h<d\right\}\right),
 \qquad \tau\geq t_{\rm gate},
\label{eq:causalreserve}
\end{equation}
Equation~(\ref{eq:causalreserve}) uses earlier days only. Candidates are
validation-selected for capacity planning; no reserved job enters payment
before ledger commitment.

The cumulative state reduces explicit-prefix nonzeros from quadratic to
$O(|\mathcal S||\mathcal K||\mathcal D||\mathcal T|)$. Equation
(\ref{eq:rollingterminal}) retains observed future arrivals in the 512-slot
continuous-horizon certificate. A lexicographic pair of linear programs first
minimizes L1 distance and then operating cost on that face, so no penalty
weight trades fit against cost; all programs are solved globally with HiGHS
\cite{huangfu2018highs}. Fig.~\ref{fig:framework} summarizes the chain and
Fig.~\ref{fig:methoddetail} details the ledger-to-LP path; execution replay and
interval audits remain separate experiments.

\begin{figure}[t]
\centering
\includegraphics[width=\columnwidth]{figures/fig_method_detail.pdf}
\caption{Sparse cumulative verifier detail. Each admissible edge is a declared
release-to-deadline service opportunity; the global projection, risk fit, and
payment gate operate on the same feasible state and do not repair missing
deadlines or execution records.}
\label{fig:methoddetail}
\end{figure}

\subsection{Signed nodal and N--1 settlement}

Let $\lambda^0\in\partial V(p^0)$ be the nodal demand-value subgradient. Signed
linear settlement is $(\lambda^0)^\top(p^0-p^1)$; (\ref{eq:value}) is the
realized system-value benchmark and bilateral contract payoff, not an ISO market
payment. The convex subgradient inequality \cite{boyd2004convex} yields
\begin{equation}
 0\leq(\lambda^0)^\top(p^0-p^1)-\Delta V
 =V(p^1)-V(p^0)-(\lambda^0)^\top(p^1-p^0).
 \label{eq:bregman}
\end{equation}
Thus the linear rule upper-bounds exact value in the implemented polyhedral
model; equality requires a common supporting dual subgradient.

For market-compatible remuneration, let $\mathcal W$ contain the complete
response cycle from anticipatory shift through recovery. The space--time rule is
\begin{equation}
 \Pi^{\rm ST}=\dt\sum_{t\in\mathcal W}\sum_d
 \lambda^0_{dt}(p^0_{dt}-p^1_{dt}).
 \label{eq:spacetime}
\end{equation}
This applies nodal virtual-link remuneration \cite{zhang2020virtuallinks,zhang2022remunerating}
to the verified counterfactual; summing site payments reproduces
(\ref{eq:spacetime}).

Space--time energy remuneration is one component of the contracted DR product.
For participating sites $\mathcal D^{\rm part}$, event intervals $\mathcal E$, and
cleared price $r^{\rm DR}$, delivered capacity and operator value are
\begin{align}
 Q^{\rm cap}&=\dt\sum_{d\in\mathcal D^{\rm part}}\sum_{t\in\mathcal E}
 (p^\emptyset_{dt}-p^1_{dt}), \label{eq:capacityservice}\\
 G&=r^{\rm DR}Q^{\rm cap}+\Pi^{\rm ST}, \label{eq:operatorvalue}
\end{align}
where $p^\emptyset$ is the minimum-cost no-event schedule under the same arrival
window. Capacity credits and ex-post participation follow established DR
mechanism design \cite{satchidanandan2023twostage,chen2021incentive}; the signed
$\Pi^{\rm ST}$ term debits transfer and recovery. Equations
(\ref{eq:capacityservice})--(\ref{eq:operatorvalue}) define the operator value.

Let $x^1$ and $x^\emptyset$ induce $p^1$ and $p^\emptyset$; participant
opportunity cost is $C_{\mathrm{opp}}=\mathcal C_{\rm w}(x^1)-\mathcal C_{\rm w}(x^\emptyset)$ and surplus is
$S=G-C_{\mathrm{opp}}$. With zero disagreement utilities, the symmetric Nash
solution \cite{nash1950bargaining} activates iff $S\geq0$ and sets
\begin{equation}
 z=\mathbb I\{S\geq0\},\quad
 P^{\rm B}=z\left(C_{\mathrm{opp}}+\frac{S}{2}\right),\quad
 u_{\rm dc}=u_{\rm op}=z\frac{S}{2}.
 \label{eq:nashcontract}
\end{equation}
Thus (\ref{eq:nashcontract}) makes nonactivation explicit; activated contracts
have nonnegative utilities and an exactly balanced transfer, reported separately
from ISO nodal remuneration and exact dispatch value. The complete window
$\mathcal W$ also debits pre-event shifts and rebound, with
$\sum_{d,t\in\mathcal W}(p^0_{dt}-p^1_{dt})\dt=0$ up to tolerance.

Security-aware settlement augments (\ref{eq:sced}) with every finite
non-islanding outage. If $L_{\ell k}$ is the LODF from outage $k$ to monitored
line $\ell$, then
\begin{equation}
 -\bar f_\ell\leq
 (H_\ell+L_{\ell k}H_k)(C_gg-p)\leq\bar f_\ell,
 \quad \ell\neq k.
 \label{eq:n1}
\end{equation}
Equation (\ref{eq:n1}) is the standard preventive N--1 LODF formulation
\cite{tejada2018lodf}; all credible rows are imposed simultaneously. Islanding
outages are reported but excluded because a connected single-reference PTDF
cannot represent separate islands, consistent with the DC planning criterion
\cite{nerc2020tpl}; native ratings are retained.

\subsection{Scenario-robust N--1 payment certificate}

The pointwise envelope controls credited energy but not network payment when
different demand vectors activate different security constraints. Let
$\mathcal P=\operatorname{conv}\{p^{(1)},\ldots,p^{(L)}\}$ contain the
workload-feasible schedules and let $p^{\rm cap}$ be the contiguous-validation
single projection. The feasible-quantile profile remains an external transfer
comparator. Let $\Xi_4$ be the held-out 1st/10th/50th/90th percentiles of
measured-to-predicted GPU energy; these scenarios define the telemetry-backed
conversion contract \cite{samsi2021supercloud}. For facility profile
$p^{\rm fac}_{dt}$, fixed demand $P^{\rm fix}_d$, and benchmark scale
$s^{\rm dc}$, the nodal injection under factor $\xi$ is
\begin{equation}
 p^{\rm dc,\xi}_{dt}=s_{\rm dc}\bigl(P^{\rm fix}_d+\xi[p^{\rm fac}_{dt}-P^{\rm fix}_d]\bigr),
 \qquad p^{\rm fac}_{dt}\geq P^{\rm fix}_d .
 \label{eq:fixedflexconversion}
\end{equation}
The scale $s_{\rm dc}$ is calibrated on the held-out $q_{99}$ flexible peak;
the fixed component is invariant to $\xi$, and raw $q_{99}$ is retained without
clipping. The flexible-only factor in Experiment~16 is diagnostic. Below,
$M p^{\rm dc,\xi}_t$ is the nodal injection; the cap profile and simplex profile
use (\ref{eq:fixedflexconversion}).
For each day, \PaymentSafe{} defines
$B_\xi^{\rm cap}=\dt\sum_{t\in\mathcal E}
V_{N-1}(p_t^{\rm native}+M p_t^{\rm dc,\xi,cap})$ and solves
\begin{align}
\operatorname{lexmin}_{\alpha,e,\eta,\{g_t^\xi\}}\quad&
 \left(\sum_{d,t\in\mathcal E}e_{dt},-\eta\right)
 \label{eq:paymentcert}\\
\mathrm{s.t.}\quad&
 p^\alpha=\sum_{\ell=1}^{L}\alpha_\ell p^{(\ell)},\nonumber\\
&\alpha\geq0,\quad\mathbf1^\top\alpha=1,\quad 0\leq\eta\leq1,\nonumber\\
&e_{dt}\geq\pm(p^\alpha_{dt}-p^{\rm safe}_{dt}),\nonumber\\
&g_t^\xi\ \text{satisfies (\ref{eq:sced}) and (\ref{eq:n1})},\nonumber\\[-2pt]
&\hspace{2em}\text{with }p_t=p_t^{\rm native}+M p_t^{\rm dc,\xi,\alpha},\nonumber\\[-2pt]
&\hspace{7.5em}t\in\mathcal E,\quad \xi\in\Xi_4,\nonumber\\
&\dt\sum_{t\in\mathcal E}\sum_iC_i(g_{it}^\xi)\nonumber\\[-2pt]
&\hspace{2em}\leq(1-\eta)B_\xi^{\rm cap},
\quad \xi\in\Xi_4.\nonumber
\end{align}
The same piecewise-linear epigraph is used for settlement. Two linear programs
first minimize trajectory error and then maximize the common fractional margin
$\eta$ on that face. Feasible dispatches in (\ref{eq:paymentcert}) give the
upper bound on optimal N--1 value; unit weight on $p^{\rm cap}$ proves
nonemptiness. Thus, for any $\xi\in\Xi_4$ and realized load $p^{1,\xi}$,
daily payment obeys
\begin{multline}
\sum_{t\in\mathcal E}\dt
[V_{N-1}(p_t^{\rm native}+M p_t^{\rm dc,\xi,\alpha})
-V_{N-1}(p^{1,\xi}_t)]\\
\leq\sum_{t\in\mathcal E}\dt
[V_{N-1}(p_t^{\rm native}+M p_t^{\rm dc,\xi,cap})
-V_{N-1}(p^{1,\xi}_t)].
\label{eq:paymentdominance}
\end{multline}
The realized-cost term cancels, so the certificate needs no execution truth and
holds samplewise; every finite non-islanding contingency is included. Standard
SCED and epigraph foundations are used \cite{zimmerman2011matpower,boyd2004convex}.
$\mathcal P$ contains the six first-stage projections and the feasible-quantile
projection. \RiskSafe{} is the validation-fitted simplex point under total and
CVaR budgets; the separately stored payment-contract profile uses the
pointwise cap and is not fed back into the risk fit.

The finite cap is complemented by a set-valued payment certificate. Let
$p^{\rm cap}=p^{(\ell^\star)}$ be the contiguous-validation cap and let
$p^{\rm pay}$ be the independently payment-selected target. For
$\Delta V_\xi(p)=V_{N-1}(p,\xi)-V_{N-1}(p^{1,\xi})$ and
$p_\theta=(1-\theta)p^{\rm cap}+\theta p^{\rm cert}$, the induced interval is
 \begin{equation}
  \mathcal I_\xi=\left[
  \min_{0\leq\theta\leq1}\Delta V_\xi(p_\theta),\;
  \max\{\Delta V_\xi(p^{\rm cap}),\Delta V_\xi(p^{\rm cert})\}\right].
  \label{eq:paymentinterval}
 \end{equation}
The lower endpoint is one exact joint N--1 LP with shared segment parameter;
convexity gives the upper endpoint. The cap payment lies in $\mathcal I_\xi$.
Endpoints use validation-frozen profiles and the realized counterfactual only;
no oracle baseline selects or certifies the interval. Its seven-profile hull
is the finite-scenario certificate in (\ref{eq:paymentcert}).

\subsection{Participant allocation}

For participants $N$, let $v(S)$ be signed dispatch value when coalition $S$
moves from baseline to actual demand. The exact allocation is
\begin{equation}
 \phi_i(v)=\sum_{S\subseteq N\setminus\{i\}}
 \frac{|S|!(|N|-|S|-1)!}{|N|!}
 [v(S\cup\{i\})-v(S)].
 \label{eq:shapley}
\end{equation}
This is the Shapley value \cite{shapley1953value}; signed grid value is the
study-specific characteristic function. Grouping (\ref{eq:shapley}) by coalition
size gives $\sum_i\phi_i=v(N)$ for $v(\varnothing)=0$. All 16 four-site and 256
eight-site coalitions are enumerated exactly. Larger contracts use exchangeable
site slices with binomial multiplicity, evaluating only $\prod_g(m_g+1)$ states;
the exact value is checked through 20 participants.

\subsection{Analytical guarantees}

The certificates make the scope explicit. Under the assumptions of
(\ref{eq:profit}), a reference-day perturbation is profitable iff
$q\pi^{\rm DR}/H>\min_jc'_j$ in (\ref{eq:threshold}); Experiment~1 checks this
against the independently solved 56-cell grid. The workload set is convex, so
every simplex schedule and (\ref{eq:twosided}) are feasible. For any execution
$c$,
\begin{align*}
[p^{\rm safe}_{dt}-c_{dt}]_+&\leq[p^{(\ell^\star)}_{dt}-c_{dt}]_+,\\
[c_{dt}-p^{\rm safe}_{dt}]_+&\leq[c_{dt}-p^{\rm floor}_{dt}]_++\varepsilon .
\end{align*}
Summation gives both bounds without locked outcomes.

The base-case and stacked N--1 values are convex piecewise-affine functions of
nodal demand. Equation (\ref{eq:bregman}) gives
$0\leq\Delta V_{\rm lin}-\Delta V\leq
\lVert\lambda^1-\lambda^0\rVert_*\lVert p^1-p^0\rVert$; positive-only
settlement has no such identity. In (\ref{eq:paymentcert}), each embedded
dispatch is a feasible witness, so subtracting the common realized cost gives
(\ref{eq:paymentdominance}) for every $\xi\in\Xi_4$. The dual representation
also makes $V_{N-1}(p;\xi)$ convex in a continuous conversion factor, so the
worst case on $[\xi_{\rm L},\xi_{\rm U}]$ is attained at an endpoint. The
lower endpoint of (\ref{eq:paymentinterval}) is the exact joint-LP minimum
over the declared affine segment; convexity makes its endpoint maximum exact.
Experiments~9 and 15 use independent 40- and four-segment models.

Finally, SHA-256 makes post-commit changes to the scheduler tuple detectable in
$h_{\rm sub}$, while joins check ordering and energy conservation after the event.
An optimum of (\ref{eq:jobindexed}) enforces entitlement, submit-time windows,
site capacity, and the requested-GPU nameplate without aggregate rounding;
contiguous execution is audited separately.
With $B_{dj}=\mathbf1\{s_j=d\}$ and bus map $M$, the job, aggregation, and
mapping residuals are
\begin{align*}
r^{\rm job}_j&=\sum_{\tau=r_j}^{d_j-1}u_{j\tau}-\widehat E_j,\\
r^{\rm agg}_{d\tau}&=X_{d\tau}-\sum_jB_{dj}u_{j\tau},\\
r^{\rm map}_{b\tau}&=p_{b\tau}-P^{\rm fix}_{b\tau}-\dt^{-1}\sum_dM_{bd}X_{d\tau}.
\end{align*}
If these typed norms are at most $\varepsilon$, the profile is the same indexed
witness up to tolerance; at zero residual an independently optimized aggregate
cannot receive settlement value. Experiment~22 checks the residuals before
N--1 valuation and preserves $\sum_i\phi_i=v(N)$.

\section{Experimental Design}\label{sec:design}

\subsection{Data, preprocessing, and information boundary}

BurstGPT supplies 5,188,507 successful requests and token service
\cite{wang2024burstgpt}; MIT Supercloud supplies scheduler submissions and
independent Data Center GPU Manager (DCGM) energy. Deterministic validation
leaves 71,144 submissions and 71,128 positive-energy matches, with 68,664
starting in the common 121-day tensor \cite{samsi2021supercloud}. Submission
time defines arrivals and measured intervals define scoring truth. A held-out
GPU-power model (28,413 observations) gives $R^2=0.899$, MAE 11.22 W, and RMSE
15.82 W; scaling uses the first 40 complete training days. Held-out ratios
define $\Xi_4$ in (\ref{eq:paymentcert}); Experiment~14 retains all matches,
whereas Experiment~19 uses declarations and a training-only entitlement. The
execution join is opened only for post-event scoring. Its runtime window
plus the precommitted 96-slot queue allowance is a contract-window definition,
not a submit-to-completion interpretation of Slurm ``timelimit''.

Because utility bus locations are unavailable, four-region coupling is a
trace-driven benchmark: a feature-stratified round robin assigns class, GPU
count, runtime, and submission time. Experiment~11 exhausts 24 permutations,
three penetrations, and two concentration controls; Experiment~18 uses
fixed generator-bus role strata at native-load multiplier 0.90 and is therefore
conditional on the declared trace-to-bus scenario. Held-out ratios span
0.5845--2.6931. Equation (\ref{eq:fixedflexconversion}) carries the fixed
6-MW/site component separately and calibrates on $q_{99}$, leaving both raw
endpoints activation-eligible; the 0.8901 flexible-only factor is planning
diagnostic. PGLib/MATPOWER/PYPOWER provide topology, ratings, and costs
\cite{babaeinejadsarookolaee2021pglib,zimmerman2011matpower}; RTS-24 follows
the published case \cite{grigg1999rts}.

\subsection{Baselines, parameters, and inference}

Experiment~2 compares eight meter/statistical baselines with complete-ledger
quantile, feasible, synthetic-control, \TailRisk, single-projection, and
\RiskSafe{} profiles; \PaymentSafe{} is evaluated separately in Experiments~9
and 15 because its hull and N--1 objective differ. Feature equality is audited
with execution truth withheld. Deferral windows are 0, 4, and 512 slots; each
region has a 118-MW flexible nameplate and 6-MW fixed component. Six projection
penalties, four risk reserves, and four response weights are validation-only.
The gate commits 5 MWh, payment uses an independent 40-segment evaluator, and
three-day bootstrap/sign tests use Holm correction
\cite{kunsch1989bootstrap,holm1979multiple}. Sparse-system size grows from
5,760 variables/13,424 nonzeros at 96 slots to 36,480/85,488 at 608 slots
(0.039--0.212 s).

The model-out check uses polar alternating-current optimal power flow (AC OPF)
with active/reactive, voltage, and apparent-power limits
\cite{zimmerman2011matpower}. After the intact solve, every outage imposes
for each non-reference generator $i$,
\begin{equation}
 P_{Gi}^{(k)}=P_{Gi}^{(0)},\qquad i\notin\mathcal G_{\rm ref},\ k\in\mathcal C.
 \label{eq:acplan}
\end{equation}
The reference generator balances active losses; $Q_G^{(k)}$ and voltages remain
recourse. Limits are checked after each fixed-plan solve, so
(\ref{eq:acplan}) tests one shared active plan. Non-reference bounds are fixed
at intact output $\pm10^{-8}$ MW with a $10^{-6}$-MW residual margin; the largest
deviation is $2.52\times10^{-8}$ MW.

Experiments~1--2 test incentives and workload verification; 3--12 test signed
settlement, transfer, risk, N--1/AC value, and continuous-cycle remuneration;
13--16 provide measured replay, indexed witness, conversion intervals, and
ledger reconciliation. Experiments~17--25 freeze the gate, verify AC and
job/network coupling, audit scale, replay an independent event, and enumerate
RTS-24 outages. Confirmatory comparisons use the same 54 locked days unless a
deterministic network or ledger rule is stated.

\section{Results}\label{sec:results}

\subsection{Incentive boundary and independent verification}

Experiment~1 solves the complete $8\times7$ DR-price/call-probability grid.
The dual threshold (\ref{eq:threshold}) matches the strategic program in all 56
cells. At \$150/MWh and $q=0.35$, reference scheduling inflates the event
baseline by 8.223 MWh; 44.458 of 55.452 MWh credited is unsupported (ratio
0.802). The profitable region follows the derived marginal boundary without a
fitted classifier.

Experiment~2 evaluates twelve methods on 54 later days (Table~\ref{tab:base}).
The CVaR-only ablation changes (nRMSE, $F_1$, false credit) from
$(0.339,0.418,1.823)$ to $(0.296,0.518,2.521)$ MWh/day, while the joint
total-plus-CVaR fit gives $(0.314,0.470,1.809)$. The latter is \RiskSafe{};
its payment cap is stored separately. Against feasible quantile, false credit
falls 2.811 to 1.809 MWh/day while nRMSE rises 0.301 to 0.314, exposing the
risk--accuracy frontier. Relative to the single reference, total-plus-CVaR
changes (nRMSE, false credit, $F_1$) by $(-0.025,-0.014,+0.051)$ and CVaR-only
by $(-0.043,+0.698,+0.100)$; removing both constraints gives 3.241 MWh/day.

\begin{table}[t]
\centering
\caption{Locked 54-day mechanism-isolation response results (daily means).}
\label{tab:base}
\footnotesize
\resizebox{\columnwidth}{!}{%
\begin{tabular}{lrrrr}
\toprule
Method & nRMSE & False MWh & Under-credit & $F_1$\\
\midrule
High-5-of-10 & 4.196 & 172.455 & 6.474 & 0.143\\
Ridge & 2.761 & 109.817 & 5.306 & 0.233\\
Gradient Boosting & 2.343 & 96.077 & 4.817 & 0.245\\
Extra Trees & 3.350 & 133.961 & 6.804 & 0.171\\
Metadata Gradient Boosting & 2.201 & 92.874 & 4.193 & 0.259\\
Ex-post Metadata Gradient Boosting & 2.126 & 90.807 & 4.184 & 0.265\\
Ex-post Quantile Gradient Boosting & 0.932 & 23.144 & 6.085 & 0.533\\
Synthetic Control & 2.598 & 104.998 & 7.877 & 0.202\\
Feasible Quantile Projection & 0.301 & 2.811 & 10.772 & 0.519\\
Single Feasible Projection & 0.339 & 1.823 & 13.537 & 0.418\\
Tail-Risk Feasible Counterfactual & 0.296 & 2.521 & 10.864 & 0.518\\
\textbf{Risk-Constrained Convex Verifier} & \textbf{0.314} & \textbf{1.809} & \textbf{12.517} & \textbf{0.470}\\
\bottomrule
\end{tabular}%
}
\end{table}

The independent trace-meter replay is a deployment check: event-window MAE is
2.673/2.651/2.810/2.988 MW for feasible-quantile, \TailRisk, \RiskSafe, and
single-projection profiles on the same 54 days. It measures closed-DCGM
alignment, not the mechanism-isolation panel.

Closest domain baselines use identical arrivals, deadlines, capacities, event
slots, and locked days. The incentive-compatible spatial-DR analogue gives
$(\mathrm{nRMSE},\mathrm{false\ credit},F_1)=(3.068,114.776,0.137)$; the
flexibility/receding-horizon controls give $(0.339,1.823,0.418)$. Han-style
joint-capacity and Chen-style coupled-regulation analogues are structural
controls (the latter and the all-site exact comparator give $(0.326,0,0.285)$),
not software reimplementations \cite{han2026dispatchable,chen2021idccoupling};
receding-horizon translation follows \cite{zhang2023receding}. Temporal-only
and joint spatio-temporal controls have migration 0 and 0.206 MWh. The ablation
panel contains 810 global schedules (binding shares 0.310/0.948; slope 1.003).

\subsection{Nodal value and spatial response}

Experiment~3 crosses thirteen counterfactuals with five settlement mechanisms
(Table~\ref{tab:settlement}). Exact signed nodal value has MAE \$151.719 versus
\$360.580 for gross response; workload errors are \$455.865 and \$421.677.
Signed netting reduces trace error by \$208.190/day (Holm-adjusted
$p=3.89\times10^{-4}$), while uniform-to-nodal and linear-to-exact changes are
$-\$1.463$ ($p=0.0118$) and \$0.671 ($p=0.8026$); workload effects are
$-\$34.165$, $-\$1.191$, and $-\$0.023/day$. All 5,616 instances satisfy
(\ref{eq:bregman}), with maximum violation $2.03\times10^{-11}$ dollars.

\begin{table}[t]
\centering
\caption{Base-case and complete N--1 settlement errors (USD/day).}
\label{tab:settlement}
\scriptsize
\resizebox{\columnwidth}{!}{%
\begin{tabular}{llrrcc}
\toprule
Panel & Settlement & Trace MAE & Workload MAE & Trace median/overpay & Workload median/overpay\\
\midrule
Base case & Nodal gross & 360.580 & 421.677 & -- & --\\
Base case & Nodal exact net value & \textbf{151.719} & 455.865 & -- & --\\
N--1 & Base exact & 139.221 & 167.367 & 79.79/130.42 & 121.01/107.69\\
N--1 & N--1 linear & 2.208 & \textbf{150.461} & 0.60/1.91 & 113.43/20.38\\
N--1 & \textbf{N--1 exact} & \textbf{0.990} & 150.955 & 0.87/0.01 & 113.43/20.18\\
\bottomrule
\end{tabular}%
}
\end{table}

Experiment 4 applies the predeclared largest gross-minus-net offset (day 87),
exposing the rebound component that positive-only local accounting omits.

\subsection{Cross-network and N--1 security evidence}

Experiment~5 evaluates 6,480 four-network outcomes; exact value improves all
12 trace cells (mean \$1.913/day; eight Holm-adjusted tests below 0.05) with a
\$0.050/day workload change. Experiment~6 solves 810 capacity/deadline schedules
(binding shares 0.031--0.911; false credit 2.374--2.740 MWh). Experiment~7
enumerates 16 four-site, 256 eight-site coalitions, and 1,296 exact
20-participant states (maximum residual $5.68\times10^{-14}$ dollars).
Experiment~8 imposes (\ref{eq:n1}) for all 37 RTS-24 non-islanding outages:
trace/workload MAE is \$139.221/\$2.208/\$0.990 and
\$167.367/\$150.461/\$150.955 for base exact/N--1 linear/N--1 exact, with
1.000 p.u. maximum loading.

Experiment~9 tests (\ref{eq:paymentdominance}) with 54 global programs and a
validation-frozen seven-profile hull. Contiguous-fold baseline nRMSE selects the
cap, while independent validation N--1 payment MAE selects the target; locked
outcomes are unopened. Held-out factors 0.584/0.734/0.996/1.278/2.693 have cap
violations below $3\times10^{-9}$ USD. Below q99, independent 40-segment MAE is
\$70.792/\$59.309/\$62.069/\$60.073 per day for \PaymentSafe, single,
feasible-quantile, and \RiskSafe; mean overpayments are
\$2.947/\$7.727/\$6.954/\$7.502. At q99, MAE is
\$214.314/\$179.816/\$188.315/\$182.297 for the same methods, with the cap
proof unchanged. Experiment~15 retains raw $q_{01}=0.5845$ and $q_{99}=2.6931$;
fixed/flexible conversion (\ref{eq:fixedflexconversion}) makes both endpoints
activation-eligible under q99 calibration without clipping. The flexible-only
0.8901 factor is diagnostic. Unseen 0.80/1.20 transfer gives certified MAE
\$16.170/\$24.262 and feasible-quantile MAE \$17.108/\$25.678. The conversion,
preventive-AC, and scale panels use the same resource semantics in
Table~\ref{tab:security}.

Table~\ref{tab:security} separates the preventive DC certificate, the
shared-active-plan AC envelope, and the all-outage AC replay.

\begin{table}[t]
\centering
\caption{Cross-network security panels and declared recourse scope.}
\label{tab:security}
\scriptsize
\begin{tabular}{@{}p{0.36\columnwidth}p{0.24\columnwidth}p{0.28\columnwidth}@{}}
\toprule
Panel & Evaluated cells & Maximum loading / status\\
\midrule
Preventive DC N--1 (RTS-24) & 37 outages; 864 cells & 1.0000 p.u.; all feasible\\
Shared-active-plan AC (IEEE-9) & 3,888 cells & 0.880/0.907/0.935 p.u. at 3/6/9\%\\
Corrective AC replay (IEEE-9/14) & 864 base + 5,400 outage cells & all connected cells feasible\\
All-outage AC replay (RTS-24/30/39/118) & 864 cells; 31,968 solves & 1.0000 p.u.; all feasible\\
\bottomrule
\end{tabular}
\end{table}

\begin{figure*}[t]
\centering
\includegraphics[width=0.92\textwidth]{figures/fig17_cross_layer_robustness.pdf}
\caption{Cross-layer robustness evidence. (a) The four-factor payment
certificate remains below its reference cost. (b) The shared-active-plan AC
envelope stays below the unit loading bound at all declared penetrations. (c)
Spatial transfer error is reported across the complete 24-mapping panel and
three penetrations.}
\label{fig:crosslayer}
\end{figure*}

Figure~\ref{fig:crosslayer} visualizes the same three evidence layers.

Experiment~10 replaces the linear network with AC OPF
\cite{zimmerman2011matpower}: 864 base and 5,400 corrective IEEE-9/14 outage
outcomes satisfy the limits. Experiment~18 separately fixes the intact active
plan; its 3,888 cells have zero active deviation and maximum loading
0.880/0.907/0.935 p.u. at 3\%/6\%/9\%. Experiment~11 spans 24 permutations,
two concentration controls, three penetrations, four counterfactuals, and
54 days; averaged over 26 assignments, \RiskSafe{} MAE is
\$98.839/\$198.726/\$296.240 at 3\%/6\%/9\%, versus
\$105.250/\$211.555/\$315.490 for single and
\$81.613/\$164.249/\$244.481 for feasible quantile. Experiment~16 audits
66,769/28,413 training/test jobs (MAE 11.223 W, RMSE 15.818 W, $R^2=0.8994$).
Experiment~12 reports event-window MAE
\$26.937/\$15.969/\$5.325/\$5.325 and complete-cycle residual
\$0.621/\$0.607/\$0.000/\$0.000 for feasible-quantile,
payment-certified, RiskSafe, and single, respectively; the zero residual is an
accounting certificate on the independent optimum face. All 54 contracts
activate.

\subsection{Independent trace and deployment-boundary audits}

Immutable MIT scheduler/DCGM replay gives slot MAE 3.740 versus 5.678 MW at
100\% coverage and a 71,128-job residual below $10^{-17}$ MWh. Observational
trace-meter MAE is 2.988 MW for \RiskSafe{} (2.673 for feasible quantile);
no utility event label is available. Experiment~23 applies a predeclared tariff
and independent LP to a submission-masked ledger with no target, risk envelope,
hull, or verifier output. Across 54 days the
gate-committed response gives nRMSE/false credit/payable/precision/recall/$F_1$
=2.606/145.506/227.116/0.351/0.436/0.358, while the independent policy gives
1.710/29.952/38.378/0.188/0.038/0.054 (Table~\ref{tab:decision}). This is a
structural non-reuse probe, not a superiority or field-treatment comparison.
Validation selects \$10/MWh, weight 1.0, and a 20-MWh/day budget; reserve
$\eta=0.60$ is planning-only. Experiment~18 supplies 9,648 fixed-plan AC-OPF
outcomes at multiplier 0.90 with zero voltage violation and maximum loading
1.000000 p.u.

\begin{table}[t]
\centering
\caption{Information-boundary and locked response certificates (54-day means). The independent policies use a submission-masked ledger, distinct tariffs, and no verifier output; the rows therefore test contract/response separation rather than a field treatment effect.}
\label{tab:decision}
\footnotesize
\resizebox{\columnwidth}{!}{%
\begin{tabular}{lrrrrrr}
\toprule
Method & nRMSE & Gross false MWh & Payable MWh & Meter underpay & Recall & F1\\
\midrule
Independent gate-committed response & 2.606 & 145.506 & 227.116 & 94.055 & 0.436 & 0.358\\
Independent risk-constrained response & 1.710 & 29.952 & 38.378 & 167.240 & 0.038 & 0.054\\
Complete-ledger \RiskSafe{} (mechanism isolation) & 0.339 & 1.823 & 10.009 & 13.537 & 0.335 & 0.418\\
\bottomrule
\end{tabular}%
}
\end{table}

\subsection{Indexed workload counterfactual}

Experiment~19 assigns each submitted job and admissible slot a variable. Its
submit-time runtime-plus-96-slot window, 0.001-MW/GPU bound, declared region,
entitlement, and 118-MW capacity yield 12,296,675 variables through an exact
separable knapsack (Table~\ref{tab:jobcounterfactual}); event energy is
0.194/0 MWh with zero residual and 29.486 MWh slack. Experiment~22 checks the
typed invariant over 1,056 event slots and all finite RTS-24 outages: raw secure
value is \$0.004293 per interval without reoptimization. The identical indexed
vector gives \$0.006352 at 1.479689$\times$ fixed-nameplate and
\$9.213788 at 2146.204351$\times$ capacity-proportional scale; all replays
pass 37 N--1 checks. Experiment~21 reports scale factors, not a second
trajectory.

Experiment~25 independently validates the declaration boundary: 71,128/71,144
jobs match execution, requested-nameplate coverage is 1.000, and no window is
infeasible under the committed rule. Its 689-job stress cohort uses 0.001 MW
per region and a 0.004-MWh floor; 43.75\% of capacity slots bind and minimum
slack exceeds $-10^{-19}$ MWh. This is separate from Experiment~19
(Table~\ref{tab:jobcounterfactual}).

Experiment~24 freezes two profiles and evaluates 37 finite RTS-24 outages over
54 days and 8 event slots (864 cells; 31,968 solves). Every cell succeeds with
maximum loading 1.000000 p.u.; the outage count and bound are in
Table~\ref{tab:jobcounterfactual}.

\begin{table}[t]
\centering
\caption{Exact indexed job-level counterfactual and binding-capacity certificates.}
\label{tab:jobcounterfactual}
\scriptsize
\begin{tabular}{@{}p{0.58\columnwidth}p{0.32\columnwidth}@{}}
\toprule
Quantity & Value\\
\midrule
Submitted / execution-matched jobs; job--slot variables & 71,144 / 71,128; 12,296,675\\
Native / counterfactual event energy; gross/net; rebound (MWh) & 0.194 / 0; \textbf{0.194}/\textbf{0.194}; 0\\
Maximum job residual; minimum site slack (MWh) & 0; 29.486\\
Binding stress cohort / capacity / service floor (jobs/MW/MWh) & 689 / 0.001 / 0.004\\
Capacity-proportional / fixed-nameplate scale ($\times$) & 2146.204 / 1.480\\
Job-to-network residual; coupled event slots (MWh / count) & 0; 1,056\\
Secure N--1 value, raw / fixed / capacity scale (USD/interval) & \textbf{0.004293} / 0.006352 / 9.213788\\
RTS-24 outages / cells / solves per coupled replay & 37 / 864 / 31,968\\
Maximum counterfactual post-contingency loading (p.u.) & 1.0000\\
\bottomrule
\end{tabular}
\end{table}

\section{Conclusion}\label{sec:conclusion}

This paper couples a provenance-bound workload verifier to signed nodal
settlement. The validation-fitted total-plus-CVaR risk profile reaches
nRMSE/false credit/$F_1=0.314/1.809/0.470$ MWh/day on the complete-ledger
mechanism-isolation panel, while the slot-60 gate replay uses only committed
declarations and is reported separately (Table~\ref{tab:decision}). The
submitted-job witness spans 12,296,675 variables, yields 0.194 MWh net
reduction with zero rebound, and maps to a zero-residual RTS-24 N--1
certificate at the raw, fixed-nameplate, and capacity-proportional scales. A
689-job stress cohort binds regional capacity at 0.001 MW while meeting the
0.004-MWh floor. Fixed facility demand is kept outside the conversion factor,
so raw $q_{99}$ is an auditable activation endpoint; the flexible-only
capacity factor remains a planning diagnostic. Deployment still requires
signed co-located telemetry and topology access, and causal utility validation
remains external.

\let\savedthebibliography\thebibliography
\let\savedendthebibliography\endthebibliography
\renewenvironment{thebibliography}[1]{\savedthebibliography{#1}\fontsize{7pt}{8pt}\selectfont}{\savedendthebibliography}
\bibliographystyle{ieeetr}
\bibliography{references}

\end{document}
